\section{Theoretical background}
\label{sec:theoretical_background}


Let $\mathcal{V} \subset \mathbb{R}^{3}$ be a \textit{regular region} with boundary
$\partial \mathcal{V}$. 
Some important properties of $\mathcal{V}$ in this case is that it is an open and 
connected set of points and $\partial \mathcal{V}$ is a closed and compact surface 
free of double points \citep[][ p. 22]{freeden_gerhards2012}.
This definition is slightly different from that given by \citet[][ p. 100]{kellogg1967},
for example, which includes the boundary as part of the region.

Consider a set of 3D gravity or magnetic sources located in the subsurface,
over $\overline{\mathcal{V}} = \mathcal{V} \, \cup \, \partial \mathcal{V}$.
Consider also that the density (contrast) or total-magnetization intensity within 
the sources can be represented by a piecewise continuous, integrable and bounded function
$\sigma(\mathbf{r}')$ defined at points 
$\mathbf{r}' = (x', y', z')^{\top} \in \overline{\mathcal{V}}$.
Here, we limit ourselves to magnetic sources having constant total-magnetization direction
defined by an unit vector $\hat{\mathbf{h}}$.

In this case, the gravitational and magnetic fields produced by the 3D sources 
at an observation point $\mathbf{r} = (x, y, z)^{\top}$ 
can be defined in terms of the following Newtonian potential:
\begin{equation}
	\Theta(\mathbf{r}) = \kappa \iiint\limits_{\mathcal{V}}
	\sigma(\mathbf{r}') \: \frac{1}{\| \mathbf{r} - \mathbf{r}'\|}
	\: dx'dy'dz' \: ,
	\label{eq:scalar-potential-Theta}
\end{equation}
where $\kappa$ is a proportionality constant; 
$\tfrac{1}{\| \mathbf{r} - \mathbf{r}' \|}$ is the inverse distance function
between the observation point $\mathbf{r}$ and a point $\mathbf{r}' \in \mathcal{V}$.
If we consider gravity sources, then $\Theta(\mathbf{r})$ 
(equation \ref{eq:scalar-potential-Theta}) represents the gravitational potential.
If we consider magnetic sources instead, then the negative directional 
derivative of $\Theta(\mathbf{r})$ (equation \ref{eq:scalar-potential-Theta}) along
the constant total-magnetization direction $\hat{\mathbf{h}}$ represents the magnetic 
scalar potential.

It can be shown that, at all points $\mathbf{r}$ located outside the sources (i.e., 
$\mathbf{r} \notin \overline{\mathcal{V}}$), $\Theta(\mathbf{r})$ (equation
\ref{eq:scalar-potential-Theta}) has continuous partial derivatives of all orders
and satisfies the Laplace's equation \citep[][ p. 122--124]{kellogg1967}.
Besides, it is regular at infinity \citep[][ p. 217]{kellogg1967}.
As a consequence of these properties, the potential $\Theta(\mathbf{r})$ also satisfies 
the following integrals:
\begin{equation}
\Theta(\mathbf{r}) =  
\int\limits_{-\infty}^{\infty} \int\limits_{-\infty}^{\infty}
\Theta(\mathbf{r}'') \:\: 
\partial_{z} \Psi(\mathbf{r}, \mathbf{r}'') \: dx''dy'' \: , \quad z < z''_{0} \: ,
\label{eq:Theta-Dirichlet}
\end{equation}
and 
\begin{equation}
\Theta(\mathbf{r}) =  
\int\limits_{-\infty}^{\infty} \int\limits_{-\infty}^{\infty}
\partial_{z} \Theta(\mathbf{r}'') \:\: 
\Psi(\mathbf{r}, \mathbf{r}'') \: dx''dy'' \: , \quad z < z''_{0} \: ,
\label{eq:Theta-Neumann}
\end{equation}
where 
$\partial_{z} \Theta(\mathbf{r}'') \equiv \tfrac{\partial \Theta(\mathbf{r}'')}{\partial z}$
and 
$\partial_{z} \Psi(\mathbf{r}, \mathbf{r}'') \equiv \tfrac{\partial
\Psi(\mathbf{r}, \mathbf{r}'')}{\partial z}$, with 
\begin{equation}
\Psi(\mathbf{r}, \mathbf{r}'') = 
\frac{1}{2\pi} \: 
\frac{1}{\| \mathbf{r} - \mathbf{r}'' \|}
\label{eq:Psi}
\end{equation}
representing the scaled inverse distance function between an observation point
$\mathbf{r} = (x, y, z)^{\top} \notin \overline{\mathcal{V}}$ and a point 
$\mathbf{r}'' = (x'', y'', z''_{0})^{\top} \notin \overline{\mathcal{V}}$ on the horizontal
plane $z''_{0}$ located above the sources.
Equation \ref{eq:Theta-Dirichlet} states that $\Theta(\mathbf{r})$ can be defined
in terms of its values $\Theta(\mathbf{r}'')$ on the plane $z''_{0}$ and 
represents the solution of the \textit{Dirichlet’s problem} or the \textit{first boundary
value problem of potential theory} \cite[][ p. 236]{kellogg1967} on a plane.
On the other hand, equation \ref{eq:Theta-Neumann} shows that $\Theta(\mathbf{r})$ can 
be defined in terms of its vertical derivative values $\partial_{z} \Theta(\mathbf{r}'')$ 
on the plane $z''_{0}$ and represents the solution of the 
\textit{Neumann’s problem} or the \textit{second boundary value problem of potential
theory} \cite[][ p. 246]{kellogg1967} on a plane.
The common outline to deduce equations \ref{eq:Theta-Dirichlet} and \ref{eq:Theta-Neumann}
can be found, for example, in \citet{roy1962} and \citet{reis_etal2020}.

